\chapter{Keyboard shortcuts and mouse controls}
\label{app:shortcuts}

Read this appendix when you already know what you want to do and only need the
keystroke for it, or when a key that is described somewhere else does not appear
to work. Chapter~\ref{ch:interface} describes the interface these keys drive,
and the commands behind them are documented in Appendix~\ref{app:console}.

Every keyboard shortcut in MeerK40t is a console command (Chapter~\ref{ch:console}).
The keymap is data, not code: the ``nudge the selection 1\unit{mm} right''
instruction that \key{Right} sends is the same \cmd{translate 1mm 0} you could
type by hand. That is why shortcuts can be moved to other keys, and why a key
with no shortcut can be given one.

Shortcuts are edited in the \panel{Keymap} window
(\menu{Edit>Settings>Key-Bindings}, or \cmd{window open Keymap} in the console).
It lists one row per binding, a \tabname{Key} and the \tabname{Command} it runs,
and applies every change at once. A command beginning with \cmd{+} runs while
the key is held down and is cancelled when it is released; all other commands
run once per press.

\begin{tipbox}[Remapping and restoring shortcuts]
\begin{enumerate}[label=\arabic*.,leftmargin=1.4em]
  \item Open \menu{Edit>Settings>Key-Bindings}.
  \item Click into the left-hand field and press the key combination you want to
        use. What you press is what is stored.
  \item Type the console command into the right-hand field, for example
        \cmd{outline 1mm}.
  \item Press \btn{Add Hotkey}. A binding that already exists for that key is
        replaced.
  \item To remove bindings, select one or more rows, right-click and choose
        \btn{Remove \ldots} (or \btn{Remove N entries}).
  \item To start again: \menu{Standard>Reset Keymap to defaults} in the same
        window restores every default; the console equivalents are
        \cmd{bind default} for the keys and \cmd{alias default} for the
        held-key timers.
\end{enumerate}
Bindings are stored in your settings, so they survive a restart. The
\panel{Keymap} window lists key names in the spelling the keymap stores them in,
for example \cmd{ctrl+shift+right}; on macOS it shows \key{Cmd} where the map
stores \ui{ctrl}.
\end{tipbox}

\section{Movement and nudging}

Two different things move here: the arrow keys move \emph{selected elements},
while the held keys \key{A}, \key{D}, \key{W} and \key{S} jog the \emph{laser
head}. Nothing selected (or nothing connected), nothing moves.

\begin{center}\small
\begin{tabular}{@{}L{3.6cm}L{5.6cm}L{5.6cm}@{}}
\toprule
\textbf{Key} & \textbf{Command} & \textbf{What happens} \\
\midrule
\key{Right}, \key{Left}, \key{Up}, \key{Down} & \cmd{translate 1mm 0} and the
\cmd{-1mm 0}, \cmd{0 1mm}, \cmd{0 -1mm} companions & The selected elements move
1\unit{mm} in that direction, once per press. \\
\key{Shift+Right} and the other arrows & \cmd{translate 10mm 0} and companions &
The same move in 10\unit{mm} steps. \\
\key{Ctrl+Shift+Right} and the other arrows & \cmd{translate 0.1mm 0} and
companions & Fine positioning: 0.1\unit{mm} per press. \\
\key{Alt+Right} and the other arrows & \cmd{translate 1spx 0} and companions &
One scene pixel (\ui{spx}): the smallest step, independent of the unit in
use. \\
\key{Alt+Shift+Right} and the other arrows & \cmd{translate 10spx 0} and
companions & Ten scene pixels. \\
\midrule
\key{A}, \key{D}, \key{W}, \key{S} & \cmd{+left}, \cmd{+right}, \cmd{+up},
\cmd{+down} & Hold to jog the laser head 1\unit{mm} every 0.1\unit{s} in that
direction; releasing the key stops it. Device required. \\
\key{Numpad Left}, \key{Numpad Right}, \key{Numpad Up}, \key{Numpad Down} &
\cmd{+translate_left}, \cmd{+translate_right}, \cmd{+translate_up},
\cmd{+translate_down} & Hold to move the selection 1\unit{mm} every
0.1\unit{s}. \\
\key{Numpad *}, \key{Numpad /} & \cmd{+scale_up} / \cmd{+scale_down} & Hold to
grow or shrink the selection by 2\,\% every 0.1\unit{s}, about its own
centre. \\
\key{Numpad +}, \key{Numpad -} & \cmd{+rotate_cw} / \cmd{+rotate_ccw} & Hold to
rotate the selection 2\textdegree{} clockwise or anticlockwise every
0.1\unit{s}. \\
\bottomrule
\end{tabular}
\end{center}

\noindent The held keys are aliases rather than simple commands: each one starts
a repeating timer, and the matching release (\cmd{-left}, \cmd{-translate_left},
and so on) cancels it. \cmd{alias} in the console lists them. The \panel{Keymap}
window spells the keypad keys \cmd{numpad_left}, \cmd{numpad_right},
\cmd{numpad_up}, \cmd{numpad_down}, \cmd{numpad_multiply}, \cmd{numpad_divide},
\cmd{numpad_add} and \cmd{numpad_subtract}.

\section{Selection, editing and document keys}

\begin{center}\small
\begin{tabular}{@{}L{3.0cm}L{5.8cm}L{6.0cm}@{}}
\toprule
\textbf{Key} & \textbf{Command} & \textbf{What happens} \\
\midrule
\key{Ctrl+D} & \cmd{element copy} & Duplicates the selected elements, in place on
top of the originals. \\
\key{Ctrl+G} & \cmd{planz clear copy preprocess validate blob preopt optimize
spool} & Builds the job from the current operations and sends it to the spooler
--- the same pipeline the ribbon \btn{Start} button runs. \\
\key{Ctrl+1} \ldots{} \key{Ctrl+5} & \cmd{bind 1 move $x $y}, \cmd{bind 2 move $x
$y}, and so on & Stores the laser head's current position in the number keys:
after pressing \key{Ctrl+3}, the key \key{3} moves the head back to that point.
Device required. \\
\key{Ctrl+Shift+C} & \cmd{align bed group xy center center} & Centres the
selection, treated as one group, on the bed. \\
\key{Delete}, \key{Numpad Del} (\cmd{numpad_delete}) & \cmd{tree selected delete}
& Deletes the selected nodes. A selection made in the scene is mirrored in the
\panel{Tree} pane, so this deletes the selected elements. \\
\key{Escape} & --- & Not a keymap entry: it is handled by the active tool, which
returns you to the selection tool. The old \cmd{pause} binding on this key has
been removed; use the \key{Pause} key instead. \\
\bottomrule
\end{tabular}
\end{center}

\noindent Selection and clipboard keys you may expect here --- \key{Ctrl+A}
(select all), \key{Ctrl+I} (invert), \key{Ctrl+C}, \key{Ctrl+V}, \key{Ctrl+X},
\key{Ctrl+Z} and \key{Ctrl+Shift+Z} (undo, redo), \key{Ctrl+O} (open project) ---
belong to the menus rather than the keymap (\S\ref{sec:shortcuts-menu}). They
work, but they do not appear in the \panel{Keymap} window and cannot be
reassigned there.

\section{Machine, windows and panels}

\begin{center}\small
\begin{tabular}{@{}L{3.4cm}L{4.8cm}L{6.6cm}@{}}
\toprule
\textbf{Key} & \textbf{Command} & \textbf{What happens} \\
\midrule
\key{Home} & \cmd{home} & Homes the laser (sends it to the endstops). Device
required. \\
\key{Pause} & \cmd{pause} & Real-time pause and resume of the running job.
Device required. \\
\key{L}, \key{U} & \cmd{lock} / \cmd{unlock} & Locks or unlocks the rail, so the
gantry cannot be pushed around by hand. Needs a driver that supports a rail
lock. \\
\key{F5} & \cmd{refresh} & Refreshes the main window: caches are invalidated and
the scene is redrawn. \\
\midrule
\key{Ctrl+F6} & \cmd{page home} & Switches the ribbon to its \tabname{Project}
page. \\
\key{Ctrl+F7} & \cmd{page design} & Switches the ribbon to a page named
\ui{design}. The stock ribbon has only \tabname{Project}, \tabname{Modify} and
\tabname{Config} pages, so on a default profile this command finds nothing. \\
\key{Ctrl+F8} & \cmd{page modify} & Switches the ribbon to its \tabname{Modify}
page. \\
\key{Ctrl+F9} & \cmd{page config} & Switches the ribbon to its \tabname{Config}
page. \\
\midrule
\key{F1} & (Help menu) & Opens the wiki page of the control under the pointer,
falling back to the general GUI page. Menu accelerator, not remappable. \\
\key{F11} & (Help menu) & The same handler as \key{F1}, from the \menu{Help}
menu. Not remappable. \\
\bottomrule
\end{tabular}
\end{center}

\noindent The F-keys that older instructions associate with windows ---
\key{F4}, \key{F6}, \key{F7}, \key{F9} and \key{F12} --- have no default binding
in this version; see \S\ref{sec:shortcuts-removed}.

\section{Test, rotary, wizard and reset bindings}

\begin{center}\small
\begin{tabular}{@{}L{3.8cm}L{5.0cm}L{6.0cm}@{}}
\toprule
\textbf{Key} & \textbf{Command} & \textbf{What happens} \\
\midrule
\key{Ctrl+E} & \cmd{circle 0.5in 0.5in 0.5in stroke red classify} & Draws a test
circle --- a radius of 0.5\unit{in} about a centre 0.5\unit{in} from the origin
on both axes, stroked red --- and hands it to classification, which assigns it
to a matching operation. \\
\key{Ctrl+R} & \cmd{rect 0 0 1in 1in stroke red classify} & Draws a
1\unit{in} \(\times\) 1\unit{in} test square with its top left corner on the
origin, stroked red, and classifies it the same way. \\
\key{Ctrl+Shift+O} & \cmd{outline 1mm} & Outlines the selection with a
1\unit{mm} offset. \\
\key{Ctrl+Shift+H} & \cmd{scale -1 1} & Mirrors the selection horizontally,
about its own centre. \\
\key{Ctrl+Shift+V} & \cmd{scale 1 -1} & Mirrors the selection vertically. \\
\midrule
\key{Ctrl+Alt+D} & \cmd{image wizard Gold} & Runs the \ui{Gold} image wizard on
the selected image. \\
\key{Ctrl+Alt+E} & \cmd{image wizard Simple} & \ui{Simple} wizard. \\
\key{Ctrl+Alt+N} & \cmd{image wizard Newsy} & \ui{Newsy} wizard. \\
\key{Ctrl+Alt+O} & \cmd{image wizard Stipo} & \ui{Stipo} wizard. \\
\key{Ctrl+Alt+X} & \cmd{image wizard Xin} & \ui{Xin} wizard. \\
\key{Ctrl+Alt+Y} & \cmd{image wizard Gravy} & \ui{Gravy} wizard. \\
\midrule
\key{Ctrl+Alt+Shift+A} & \cmd{arm;startjob} & Arms the laser and starts the job
in one press: the two steps of \btn{Arm} followed by \btn{Start}. \\
\key{Ctrl+Alt+Shift+Home} & \cmd{bind default;alias default} & Restores every
default binding and alias. \\
\bottomrule
\end{tabular}
\end{center}

\noindent A wizard works on an image, so select the image first; a locked image
is skipped and reported in the console.

\section{Keys with no default binding}
\label{sec:shortcuts-removed}

Several keys that older MeerK40t releases, and some third-party instructions,
bind are deliberately unbound in this version. The defaults live in
\file{meerk40t/core/bindalias.py} as a list of values per key, newest first; a
leading empty value means ``no binding'', and it also removes that key from a
profile that still carries an older default. The table lists the command each
key used to run, so that you can put it back with \btn{Add Hotkey} if you want
it.

\begin{center}\small
\begin{tabular}{@{}L{4.4cm}L{4.8cm}L{5.6cm}@{}}
\toprule
\textbf{Key} & \textbf{Command it used to run} & \textbf{How to get the same
thing now} \\
\midrule
\key{F4} & \cmd{window open CameraInterface} & \menu{Tools>CameraInterface};
requires a camera. \\
\key{F6} & \cmd{window open JobSpooler} & \menu{Tools>Spooler}, or the
\panel{Spooler} pane. \\
\key{F7} & \cmd{window open Controller} & \menu{Tools>Controller}. The window is
provided by the driver of the loaded device, so a device must be connected ---
with no device there is no Controller window to open. \\
\key{F8} & \cmd{dialog_path} & The path entry dialog, from the tree's
right-click menu. \\
\key{F9} & \cmd{dialog_transform} & The transform dialog, from the same
menu. \\
\key{F12} & \cmd{window open Console} & \menu{Tools>Console}, the \btn{Console}
ribbon button, or the \panel{Console} pane. \\
\key{Escape} & \cmd{pause} & \key{Pause}. \key{Escape} still cancels the active
tool. \\
\midrule
\key{Ctrl+A}, \key{Ctrl+C}, \key{Ctrl+V}, \key{Ctrl+X}, \key{Ctrl+I},
\key{Ctrl+O}, \key{Ctrl+S}, \key{Ctrl+Z}, \key{Ctrl+Shift+Z} & \cmd{element*
select}, \cmd{clipboard copy}, \cmd{clipboard paste}, \cmd{clipboard cut},
\cmd{element* select^}, \cmd{outline}, \cmd{dialog_stroke}, \cmd{undo},
\cmd{redo} & The menus now own these keys --- see
\S\ref{sec:shortcuts-menu}. The commands themselves still exist. \\
\key{Ctrl+F}, \key{Ctrl+S} & \cmd{dialog_fill} / \cmd{dialog_stroke} & The fill
and stroke entry dialogs, from the tree's right-click menu. \\
\key{Alt+C}, \key{Alt+E}, \key{Alt+R} & \cmd{cut} / \cmd{engrave} / \cmd{raster}
& Groups the selection into that kind of operation; use the \panel{Operations}
pane or the right-click menu. \\
\key{Alt+F}, \key{Alt+S}, \key{Alt+H}, \key{Alt+P}, \key{Alt+T} & \cmd{dialog_fill},
\cmd{dialog_stroke}, \cmd{dialog_path}, \cmd{dialog_flip}, \cmd{dialog_transform}
& The same entry and flip dialogs. \\
\midrule
\key{Alt+F3}, \key{Ctrl+F3} & \cmd{rotaryscale} / \cmd{rotaryview} & Both
commands still exist: \cmd{rotaryscale} scales the selected elements for rotary
work, \cmd{rotaryview} toggles the stretched scene preview. Bind them, or type
them in the console. \\
\key{Ctrl+Alt+G}, \key{Ctrl+Alt+S} & \cmd{image wizard Gold} /
\cmd{image wizard Stipo} & Duplicates of \key{Ctrl+Alt+D} and \key{Ctrl+Alt+O}
above. \\
\key{Ctrl+Alt+Shift+Escape} & \cmd{reset_bind_alias} & Reset the shortcuts in
the \panel{Keymap} window, or with \cmd{bind default;alias default}. \\
\key{Ctrl+Shift+L} & \cmd{signal lock_helper} & Removed upstream, with the
comment that it does not work. Rail lock and unlock are on \key{L} and
\key{U}. \\
\bottomrule
\end{tabular}
\end{center}

\section{Menu accelerators that are not in the keymap}
\label{sec:shortcuts-menu}

These keys are wired directly into the menus. They work exactly as the menu
entry does, they are printed next to the menu label, and they cannot be changed
in the \panel{Keymap} window --- rebinding one would mean editing the menu.

\begin{center}\small
\begin{tabular}{@{}L{2.6cm}L{5.4cm}L{6.8cm}@{}}
\toprule
\textbf{Key} & \textbf{Menu item} & \textbf{What happens} \\
\midrule
\key{Ctrl+N} & \menu{File>New} & Clears operations, elements and notes. \\
\key{Ctrl+O} & \menu{File>Open Project} & Opens a project; hold \key{Shift} to be
asked where to place it. \\
\key{Ctrl+S} & \menu{File>Save} & Saves the project as an SVG file, overwriting
the existing file. \\
\key{Ctrl+Shift+S} & \menu{File>Save As} & Saves into a new SVG file. \\
\midrule
\key{Ctrl+Z} & \menu{Edit>Undo} & Undoes the last action. \\
\key{Ctrl+Shift+Z} & \menu{Edit>Redo} & Reverts the last undo. \\
\key{Ctrl+X}, \key{Ctrl+C}, \key{Ctrl+V} & \menu{Edit>Cut}, \menu{Edit>Copy},
\menu{Edit>Paste} & Cut, copy and paste the selected elements. \\
\key{Ctrl+Shift+V} & \menu{Edit>Paste from system} & Pastes data copied in
another application. \\
\key{Ctrl+A} & \menu{Edit>Select all} & Selects every element on the scene. \\
\key{Ctrl+I} & \menu{Edit>Invert selection} & Inverts which elements are
selected. \\
\midrule
\key{Ctrl+-}, \key{Ctrl++} & \menu{View>Zoom Out}, \menu{View>Zoom In} & Zooms
the scene out or in. \\
\key{Ctrl+B} & \menu{View>Zoom to Bed} & Fits the machine bed in the window. \\
\key{Ctrl+Shift+B} & \menu{View>Zoom to Selected} & Fits the current selection
in the window. \\
\key{Ctrl+U} & \menu{View>GUI Appearance>Show/Hide UI-Panels} & Hides or shows
every pane except the scene. \\
\key{Ctrl+,} & \menu{Edit>Settings>Preferences} & Opens the \panel{Preferences}
window. \\
\midrule
\key{F1} & \menu{Help>Help} & Help for the control under the pointer (not on
macOS). \\
\key{F11} & \menu{Help>Online Reference} & The same help handler. \\
\bottomrule
\end{tabular}
\end{center}

\section{Mouse and trackpad}

The left button is for working on the design; the right and middle buttons ---
with a modifier where needed --- move the camera you look at it with.
Chapter~\ref{ch:interface} introduces these in context; this is the complete
list.

\begin{center}\small
\begin{tabular}{@{}L{4.8cm}L{10.0cm}@{}}
\toprule
\textbf{Input} & \textbf{What happens} \\
\midrule
Wheel up / down & Zooms in or out around the pointer. \\
Wheel with \key{Ctrl} & The other of zoom and pan; by default that means
\key{Ctrl}+wheel pans vertically. \\
Wheel left / right, or \key{Shift}+wheel & Pans horizontally. \\
Right-drag with \key{Alt} & Drag-zoom: dragging down or right zooms in, up or
left zooms out. \\
Middle-drag with \key{Ctrl} & Drag-zoom, identical to the above. \\
Right-drag with \key{Ctrl} & Drag-pan: the scene follows the pointer. \\
Middle-drag & Nothing; the event is consumed so that it does not disturb the
scene. \\
Trackpad pinch, or gesture zoom & Zooms in and out about the pointer. \\
Trackpad two-finger pan & Pans the scene. \\
\midrule
Left-drag on empty space & Rubber-band selection. Dragging \emph{right}
(enclosing the elements) selects everything the rectangle contains; dragging
\emph{left} (touching them) selects everything the rectangle touches. The
status bar names the method in use, and the rectangle is drawn in the matching
colour. \key{Shift} adds to the existing selection, \key{Ctrl} inverts the
selection state of what it catches. \\
Left click & Selects the element under the pointer. \\
Left-drag on the selection & Moves the selected elements. Press inside the
selection, hold for about half a second, then drag: the pointer changes and the
elements follow. Start dragging straight away and you get a rubber band
instead. Grid snapping and magnet lines apply while you drag. \\
Right click & The scene menu: background and grid switches, \ui{Move laser head
here}, \ui{Zoom Out}, \ui{Zoom In}, \menu{Zoom to Bed} and
\menu{Zoom to Selected}. \\
\bottomrule
\end{tabular}
\end{center}

\noindent Three preferences change this behaviour: \setting{mouse_wheel_pan}
makes the wheel pan instead of zoom, and \setting{mouse_pan_invert} and
\setting{mouse_zoom_invert} reverse the pan and zoom directions. They live in
\menu{Edit>Settings>Preferences}, \tabname{GUI}.

\section{Where these defaults come from}

The defaults are built from tables in \file{meerk40t/core/bindalias.py}:
\cmd{DEFAULT_KEYMAP} for the keys, \cmd{DEFAULT_ALIAS} for the held-key timers.
Each entry keeps its current value first and its older values after it, so a
default that has changed since you last saved your settings can be migrated to
the new value --- unless you had already bound that key yourself, in which case
your version wins.

\begin{itemize}
  \item \textbf{Nothing happens at all.} The key has no binding in this version.
        Compare it with \S\ref{sec:shortcuts-removed} --- \key{F12} for the
        console, \key{F6} for the spooler and \key{F7} for the controller window
        are the frequent cases --- and add it with \btn{Add Hotkey} if you want
        it.
  \item \textbf{Nothing happens, and the key is bound.} Most likely nothing was
        selected, or the command needs a machine. \cmd{translate}, \cmd{scale},
        \cmd{rotate}, \cmd{outline}, \cmd{element copy} and the alignment
        commands all work on the selection. \cmd{home}, \cmd{pause},
        \cmd{lock}, \cmd{unlock}, the \key{A}/\key{D}/\key{W}/\key{S} jogs and
        the \key{Ctrl+1} to \key{Ctrl+5} position keys need a connected
        device.
  \item \textbf{A window will not open.} Window commands such as
        \cmd{window open Controller} are registered by the driver of the loaded
        device. With no device loaded there is no window to show, so connect or
        add the device first (Chapter~\ref{ch:connect}).
  \item \textbf{The key types instead.} Keyboard events go to whatever has
        focus. Inside the console, a text field or a numeric field, a key is
        text --- click in the scene, or leave the field, first. While a creation
        or node-editing tool is active it consumes keys as well; press
        \key{Escape} to leave the tool.
\end{itemize}

\begin{behindbox}[Reading the current keymap]
\cmd{bind} with no arguments prints every binding as a numbered list;
\cmd{bind <key> <command>} sets one, and \cmd{alias} does the same for the
held-key timers. The \panel{Keymap} window is a front end to exactly those two
commands, which is why anything you can type in the console can be bound to a
key. Bindings are stored in your settings profile, and
\cmd{bind default;alias default} discards your changes and rebuilds the defaults
listed in this appendix.
\end{behindbox}
